\section{案例与研究方向}

\subsection{Visual Agent Evidence}

关键幻灯片展示了 evidence matrix 与 agent swarm 之间的耦合：

\begin{figure}[H]
\centering
\includegraphics[width=0.8\textwidth]{slide-02.jpg}
\caption{Slide 2 分析 agent swarm、data quality 与 governance 的关系，并附带 evidence matrix。}
\end{figure}

\subsection{Slides 与实践映射}

\begin{figure}[H]
\centering
\includegraphics[width=0.8\textwidth]{slide-03.jpg}
\caption{Slide 3 展示了 RLHF pipeline、ops playbook 与治理 checkpoint 的融合，是讲座中强调 actionability 的视觉证据。}
\end{figure}

Slide 3 把 timeline、monitoring、governance 串成矩阵，提醒团队在每个阶段都要有明确的可交付 artifact。

\begin{knowledgebox}{幻灯片到产出的映射}
在部署前，把 slide 中的 timeline、artifact 与 checkpoints 映射为 sprint backlog，可以确保每个 deliverable 都有人负责。
\end{knowledgebox}

\subsection{Research Directions}

未来研究可以关注：

\begin{itemize}
    \item 把 reward model calibration 与 policy evaluation auto-connected
    \item Multi-agent 身份下的 preference alignment
    \item Explainability layer 使 KL 约束可解释
    \item Sim-to-real research 让 reward model 适应更宽泛场景
\end{itemize}

Archit 还特别强调，需要把理论研究与实际 deployment 反馈循环起来：新的研究方向必须能生成可执行 playbook，而不仅仅是 paper-level insight。

\begin{knowledgebox}{研究方向的工程可解释性}
让研究成果可落地的关键是\textbf{工程可解释性}：解释 KL margin、reward model drift、agent dash board 指标，才能让产品经理和法律团队信任 alignment pipeline。
\end{knowledgebox}

\subsection{本章小结}

案例部分把 slide 上的视觉证据、agent swarm 结构与未来研究方向串起来，提醒我们 RLHF 还在快速演化。

